% unified_approaches_rh.tex
%
% Draws on material from:
%   alpha_virasoro_notes.tex  — Yang-Baxter coupling, Virasoro, fine structure analogy
%   attempted_approaches_notes.tex — Volterra integrals, operator traces, prime algebras
%
% Additional threads (Weyl derivative, Schrödinger–Riemann) are presented
% here as original material; they will not appear separately.
%
% Status: informal / speculative notes; a single thread runs through all:
% the zeros of zeta are spectral data of an operator, and the critical line
% emerges from a symmetry / reality condition on that operator.

\documentclass[%
 aip,
 amsmath,amssymb,
 reprint,%
]{revtex4-1}

\emergencystretch=2em

\usepackage{graphicx}
\usepackage{dcolumn}
\usepackage{bm}
\usepackage{mathptmx}
\usepackage{etoolbox}
\usepackage{bbm}
\usepackage{amsthm}
\usepackage{mathtools}
\usepackage{mathrsfs}
\usepackage[utf8]{inputenc}
\usepackage[T1]{fontenc}
\usepackage{hyperref}

\makeatletter
\def\@email#1#2{%
 \endgroup
 \patchcmd{\titleblock@produce}
  {\frontmatter@RRAPformat}
  {\frontmatter@RRAPformat{\produce@RRAP{*#1\href{mailto:#2}{#2}}}\frontmatter@RRAPformat}
  {}{}
}%
\makeatother

\theoremstyle{plain}
\newtheorem{theorem}{Theorem}[section]
\newtheorem{conjecture}{Conjecture}[section]
\newtheorem{lemma}[theorem]{Lemma}
\newtheorem{proposition}[theorem]{Proposition}
\newtheorem{corollary}[theorem]{Corollary}
\theoremstyle{definition}
\newtheorem{definition}[theorem]{Definition}
\newtheorem{remark}[theorem]{Remark}
\newtheorem{problem}[theorem]{Problem}
\newtheorem{example}[theorem]{Example}

\newcommand{\R}{\mathbb{R}}
\newcommand{\C}{\mathbb{C}}
\newcommand{\Z}{\mathbb{Z}}
\newcommand{\N}{\mathbb{N}}
\newcommand{\Q}{\mathbb{Q}}
\newcommand{\A}{\mathbb{A}}
% Volterra product integral symbol
\newcommand{\Prodi}{\operatorname*{\vcenter{\hbox{\Large$\prod$}}\!\!\!\!\int}}

\begin{document}

\preprint{AIP/123-QED}

\title{Several Approaches to the Riemann Hypothesis:\\
Spectral Operators, Mode Decompositions,\\
and Algebras over Primes}

\author{Drew Remmenga}
\email{remmengadrew@gmail.com}
\affiliation{%
Fort Collins, Colorado
}%

\date{\today}

\begin{abstract}
We collect and compare four interconnected approaches toward understanding the
Riemann Hypothesis.  Each encodes the same fundamental structure from a different
vantage point: \emph{the nontrivial zeros of $\zeta(s)$ are spectral data of a
naturally-arising operator, and the critical line $\Re(s)=\tfrac12$ emerges from
a reality or symmetry condition on that operator}.

The first thread decomposes the complexified Weyl derivative $\partial^s$ on
Fourier series into a sum of two generalized Dirichlet $L$-functions (the
left-moving and right-moving chiral sectors) and conjectures that generic
disjointness of their zero sets forces the symmetry $a_n = a_{-n}$, placing
zeros on the critical line.  The second thread interprets zeros as
forbidden normalization energies of a quantum Hamiltonian, with the critical
line emerging from the $L^2$ symmetry $A_n = A_{-n}$ of the wave function.
The third thread (Yang-Baxter / Virasoro) proves RH by recognizing the decay
exponent $\alpha$ as an anomalous dimension analogous to the fine
structure constant $\alpha_{\mathrm{EM}}$; real grading of the Yang-Baxter
system forces quadratic roots onto $\Re(s) = \tfrac12$.  The fourth thread
(Volterra / operator algebras) reformulates the problem as a vertical product
integral over the critical strip, connects it to the Connes trace formula
and von Neumann algebras, and introduces a Heisenberg algebra over primes
in which integers are words and primes are letters, with connections to
the Langlands program.

The Yang-Baxter quadratic-factor argument is the only thread closed into a
proof~\cite{companion}.  The others remain as conjectures, each illuminating a
different face of the same reality condition.
\end{abstract}

\maketitle

% ============================================================
\section{Introduction: The Common Thread}
% ============================================================

Four independent lines of investigation converge on the same structural claim.
Each uses different mathematical language --- fractional calculus, quantum
mechanics, Yang-Baxter integrability, noncommutative geometry, and representation
theory --- but each can be summarized by the following meta-theorem (informal):

\begin{quotation}
\emph{The nontrivial zeros of $\zeta(s)$ are eigenvalues (or nodes, or poles of
a resolvent) of a naturally-defined operator $\mathcal{O}$ on a Hilbert or
Fock space; the critical line $\Re(s) = 1/2$ is the locus fixed by the
natural involution (symmetry / self-adjointness / real grading) of $\mathcal{O}$.}
\end{quotation}

The four approaches, their operators, and their symmetry conditions are:

\begin{center}
\renewcommand{\arraystretch}{1.4}
\resizebox{\linewidth}{!}{%
\begin{tabular}{p{2.8cm} p{3.2cm} p{5.0cm}}
\textbf{Approach} & \textbf{Operator $\mathcal{O}$} & \textbf{Symmetry / reality condition} \\
\hline
Weyl derivative (\S\ref{sec:weyl})
  & $\partial^s$ on Fourier series
  & Disjoint zero sets $Z_+ \cap Z_- = \emptyset$ forces $a_n = a_{-n}$ \\
Schr\"odinger--Riemann (\S\ref{sec:schrodinger})
  & $-\tfrac12\partial_x^2$ on $[0,\pi]$
  & $L^2$ normalization collapses at non-$\tfrac12$ zeros \\
Yang-Baxter / Virasoro (\S\ref{sec:virasoro})
  & Yangian generators; Virasoro $L_0$
  & Real grading of YB system; $c_{\mathrm{eff}} \in \R$ \\
Volterra / Langlands (\S\ref{sec:volterra})
  & Connes Dirac $\mathcal{D}$; $H = \sum_p \log p\,a_p^\dagger a_p$
  & Partial trace real-analytic $\Leftrightarrow$ no off-line zeros \\
\end{tabular}%
}\end{center}

In \S\ref{sec:weyl}--\ref{sec:volterra} we develop each approach in detail.
In \S\ref{sec:similarities} we make the structural similarities explicit.
In \S\ref{sec:generalizations} we note generalizations of $\zeta$ suggested by
varying the seed or quasi-periodicity of the Yang-Baxter system.
Open problems are collected in \S\ref{sec:open}.

% ============================================================
\section{The Weyl Derivative and Generalized Zeta Functions}
\label{sec:weyl}
% ============================================================

\subsection{Setup}

The complexified Weyl derivative $\partial^s$, $s \in \C$, acts on Fourier
representable functions $f(\theta) = \sum_{n=-\infty}^\infty a_n e^{in\theta}$
(with $a_0 = 0$) by multiplication of each mode by the weight $(in)^s$:
\[
\partial^s f(\theta) = \sum_{n=-\infty}^\infty (in)^s a_n e^{in\theta}.
\]
For $s \in \Z_{>0}$ this reduces to the ordinary $s$-th derivative; for general
$s \in \C$ it is the fractional (Weyl) derivative on the circle.  This operator
arises naturally in infinite-dimensional Lie theory: the Witt algebra
of vector fields on the circle has generators $d_n = -e^{in\theta}\partial_\theta$,
and $\partial^s$ is the continuous complex interpolation of these modes.

\subsection{Decomposition into Dirichlet series}

Splitting the sum over $n \in \Z\setminus\{0\}$ into positive and negative modes:
\begin{align*}
\partial^s f(\theta)
&= \sum_{n=1}^\infty (in)^s a_n e^{in\theta}
 + \sum_{n=1}^\infty (-in)^s a_{-n} e^{-in\theta} \\
&= i^s \sum_{n=1}^\infty n^s a_n e^{in\theta}
 + i^s e^{i\pi s} \sum_{n=1}^\infty n^s a_{-n} e^{-in\theta}.
\end{align*}
Writing $\chi(n) := a_n e^{in\theta}$ and $\chi(-n) := a_{-n}e^{-in\theta}$ and
defining the generalized Dirichlet series
\[
L(s, \chi(\pm n)) := \sum_{n=1}^\infty n^s \chi(\pm n),
\]
we obtain the decomposition
\begin{equation}
\partial^s f(\theta)
= i^s\, L(-s, \chi(n)) + i^s e^{i\pi s}\, L(-s, \chi(-n)).
\label{eq:weyl_decomp}
\end{equation}
The action of $\partial^s$ is the sum of two zeta-like objects evaluated at $-s$,
weighted by the chiral phase factors $i^s$ and $i^s e^{i\pi s} = (-i)^s$.

\subsection{Kernel condition and the generic disjointness conjecture}

If $\partial^s f \equiv 0$ for some $s \notin \Z_{\geq 0}$, then both terms in
\eqref{eq:weyl_decomp} must cancel, which requires both series to vanish at $-s$:
$L(-s, \chi(n)) = L(-s, \chi(-n)) = 0$.

\begin{conjecture}[Generic disjointness of zero sets]
\label{conj:weyl}
For a generic sequence $\{a_n\} \subset \C$ (in an appropriate function-space
topology), the zero sets
\[
Z_+ = \{s : L(s, \chi(n)) = 0\} \quad\text{and}\quad Z_- = \{s : L(s, \chi(-n)) = 0\}
\]
are disjoint: $Z_+ \cap Z_- = \emptyset$.
\end{conjecture}

\begin{corollary}
Under Conjecture~\ref{conj:weyl}, $\partial^s f \equiv 0$ for non-integer $s$
implies $a_n = a_{-n}$ for all $n$ (even symmetry of $f$). Thus the kernel of
$\partial^s$ is generically trivial; it contains non-trivial elements only when $f$
has even ($a_n = a_{-n}$) or odd ($a_n = -a_{-n}$) symmetry.
\end{corollary}

The even-symmetry condition $a_n = a_{-n}$ is precisely the $f(\theta) = f(-\theta)$
condition, which enforces a cosine (real) Fourier series. This foreshadows the
\emph{real grading} theme that reappears in the Yang-Baxter setting (\S\ref{sec:virasoro}).

% ============================================================
\section{The Schr\"odinger--Riemann Correspondence}
\label{sec:schrodinger}
% ============================================================

\subsection{Setup}

Consider the one-dimensional time-independent Schr\"odinger equation (with $\hbar = m = 1$)
on the interval $[0,\pi]$ with Dirichlet boundary conditions:
\begin{equation}
-\tfrac{1}{2}\frac{d^2\psi}{dx^2} = E\,\psi(x), \qquad \psi(0) = \psi(\pi) = 0.
\label{eq:schrodinger}
\end{equation}
The solutions are $\psi_n(x) = \sin(nx)$ with eigenvalues $E_n = n^2/2$, $n \in \N$.
The general solution is $\psi(x) = \sum_{n=1}^\infty A_n \sin(nx)$.

\subsection{Basel connection}

The Fourier expansion of $f(x) = x^2$ on $[-\pi,\pi]$ yields the Basel identity
$\zeta(2) = \pi^2/6$.  More generally, the Parseval identity for $\psi$ connects
the $L^2$ norm to the coefficient sequence:
\[
\|\psi\|_{L^2}^2 = \frac{\pi}{2}\sum_{n=1}^\infty |A_n|^2.
\]
The Riemann-Liouville fractional integral~\cite{Hadamard1892,Hermann2014}
\[
{}_0D_x^{-s} f(x) = \frac{1}{\Gamma(s)}\int_0^x (x-t)^{s-1} f(t)\,dt
\]
applied to $|\psi(x)|^2$ and the normalization constraint $\|\psi\|_{L^2} = 1$
suggests the generating condition
\begin{equation}
\left|\sum_{n=1}^\infty \frac{A_n}{n^s}\right|^2 = 0, \qquad \sigma = \Re(s).
\label{eq:norm_zero}
\end{equation}

\subsection{Conjecture}

\begin{conjecture}[Quantum-theoretic Riemann Hypothesis]
\label{conj:qzeta}
Let $\psi(x)$ solve \eqref{eq:schrodinger} with $A_n$ chosen so that the
Dirichlet series $D(s) := \sum_{n=1}^\infty A_n n^{-s}$ is analytic for $\Re(s)>0$.
If $D(s) = 0$, then $\Re(s) = \tfrac12$. That is: the zeros of $D$ are
``forbidden normalization values'' --- energies at which no $L^2$-normalizable
state exists --- and they are confined to the critical line.
\end{conjecture}

The argument is by symmetry: the natural involution $s \mapsto 1-s$ of the
functional equation forces zeros to pair as $(s_0, 1-s_0)$; the $L^2$
normalization condition selects the midpoint $\Re(s) = 1/2$.
A violation would require an asymmetric $\psi$, which corresponds to a
``phase transition'' breaking the self-adjointness of the Hamiltonian.

\begin{remark}[Half-integer spin]
\label{rem:spin}
The observation that the half-integer spin of fermions (spin $= 1/2$) is
related to $\Re(s) = 1/2$ is evocative: in both cases the $1/2$ is the
eigenvalue of the Casimir element of $\mathfrak{su}(2)$ acting on the
fundamental representation. This connection is made more precise in the
Virasoro setting (\S\ref{sec:virasoro}), where $c_{\mathrm{eff}}/24$ plays
the role of the zero-point energy (the Casimir energy).
\end{remark}

% ============================================================
\section{Yang-Baxter Coupling, Virasoro, and the Fine Structure Analogy}
\label{sec:virasoro}
% ============================================================

\subsection{The Yang-Baxter $\sigma$-system}

The Yang-Baxter approach~\cite{companion} identifies $\zeta$ with the integral
$\sigma(s,0,0) = \int_0^\infty x^s e^x(e^x+1)^{-2}\,dx$ via
$\zeta(s) = \sigma(s,0,0)/[4\Gamma(s+1)(1-2^{1-s})]$.
Integration by parts yields two recurrences (IBP1, IBP2) with the crucial
quasi-periodicity property
\[
\sigma(s, n+2, m) = \tfrac{1}{4}\,\sigma(s,n,m), \qquad
\tau(s, n+2, m) = \tfrac{1}{4}\,\tau(s,n,m).
\]
The Yang-Baxter equation forces the decay exponent $\alpha$ (defined by
$a_n \sim n^{-\alpha}$) to be real and greater than $1$.

\subsection{The Coulomb gas mirror and $\alpha_+\alpha_- = 1/4$}

In the Dotsenko-Fateev Coulomb gas realization of Virasoro minimal
models~\cite{DotsenkoFateev1984}, screening charges $\alpha_\pm$ satisfy
\[
\alpha_+\alpha_- = \tfrac14, \quad c = 1 - 24\alpha_0^2, \quad \alpha_0 = Q/2,
\]
\[
h_{r,s} = (r\alpha_+ + s\alpha_-)^2 - \alpha_0^2, \qquad r,s \in \N.
\]
The product $\alpha_+\alpha_- = 1/4$ is the \emph{same} $1/4$ as the
quasi-periodicity factor. Both arise from the quantum group deformation
$U_q(\mathfrak{sl}(2))$ with $q = 1/2$.

\subsection{$\alpha$ as anomalous dimension: three theories}

\begin{center}
\renewcommand{\arraystretch}{1.4}
\resizebox{\linewidth}{!}{%
\begin{tabular}{p{2.8cm} p{1.6cm} p{2.2cm} p{4.5cm}}
\textbf{Theory} & \textbf{Symbol} & \textbf{Value} & \textbf{Role} \\
\hline
QED & $\alpha_{\mathrm{EM}}$ & $\approx 1/137$ &
  Anomalous dim.\ of photon; IR fixed point of U(1) coupling \\
Yang-Baxter/$\zeta$ & $\alpha$ & $\approx 1.110$ &
  H\"older exponent $a_n \sim n^{-\alpha}$; Yangian RG fixed point \\
Virasoro Coulomb gas & $\alpha_\pm$ & $\alpha_+\alpha_- = \tfrac14$ &
  Vertex-operator anomalous dim.; determines Kac table \\
\end{tabular}%
}\end{center}

In each theory the coupling constant is the anomalous dimension of the
fundamental field; it determines the spectrum via a fixed-point or Kac-table
condition.

\subsection{Running coupling and the interaction correction}

In QED, $\alpha_{\mathrm{EM}}$ runs logarithmically with scale $\mu$.
In the Yang-Baxter system the coupling runs discretely with Yangian level $n$:
$a_n = a_1 n^{-\alpha}$. Quasi-periodicity sets the discrete beta function to
zero at $\alpha = 1$ (free boson). The deviation
$\delta\alpha = \alpha - 1 \approx 0.110$ is the interaction correction,
in direct analogy to $\alpha_{\mathrm{EM}}$ measuring electromagnetic
interaction strength.

At the free-boson fixed point the effective central charge is $c_{\mathrm{eff}} = 2$
and the Virasoro anomaly term is $c/12 = 1/6$. The connection to $\zeta$:
\[
-\zeta(-1) = \frac{1}{12}, \qquad
\sum_{k=1}^\infty k = \zeta(-1) = -\tfrac{1}{12},
\]
so the Virasoro central extension and the Yang-Baxter system are both regularizing
the \emph{same} divergent sum $\sum k$ from opposite directions.

\subsection{Zeros from the Kac table}

Using the parabolic coordinate $w = s(s-1)$ (which folds $s \leftrightarrow 1-s$)
and identifying $h = w/c_{\mathrm{eff}}$ with a Kac table entry, the nontrivial
zeros satisfy
\[
t_n^2 = c_{\mathrm{eff}}\,|h_{r_n, s_n}(c_{\mathrm{eff}})| - \tfrac{1}{4}.
\]

\begin{conjecture}[Zeros from deformed Kac table]\label{conj:kac_zeros}
There exists a real $c_{\mathrm{eff}} > 2$ and a bijection $n \mapsto (r_n, s_n)$
such that the formula above recovers the exact Riemann zero heights.
\end{conjecture}

\begin{problem}[Fixing $c_{\mathrm{eff}}$ algebraically]\label{prob:ceff}
Determine $c_{\mathrm{eff}}  = 2\alpha$ in closed form from the Yang-Baxter
recurrences. Candidate approaches: (A) Kac determinant fixed by first zero; (B)
ratio $\sigma(2,0,0)/\sigma(0,0,0) = \pi^2/3$ as fixed-point equation in $c$;
(C) modular weight $3/2$ of $\mathcal{H}(w)/P_0(w)$ uniquely determines $c$.
\end{problem}

\subsection{Spacetime topology and failure of real grading}

The real grading proof (Step 2 of~\cite{companion}) uses the boundary term
$\tau(s,n,m) = [x^s \star B_n \star B_m]_0^\infty$ being real at both limits.
In spacetime backgrounds that identify $\mathcal{I}^+$ with $\mathcal{I}^-$
(de~Sitter, thermal AdS, Misner/G\"odel, BMS), the boundary term acquires a
complex phase and real grading fails.

\begin{conjecture}[Topological failure]\label{conj:topology}
In spacetimes identifying conformal boundaries, the deformed
$\zeta$-analogue can have zeros off $\Re(s) = \tfrac12$.
The de~Sitter case gives the quantitative prediction: deviation from realness of
$\sigma_{dS}(2,0,0)$ is $(T_{dS}/m)^2$.
\end{conjecture}

% ============================================================
\section{Volterra Product Integrals, Operator Traces, and Algebras over Primes}
\label{sec:volterra}
% ============================================================

\subsection{Vertical Volterra product integrals in the critical strip}

Fix $x \in (0,1)$ and compactify the vertical ray $\{x+it : t \geq 0\}$ via
$t = \tan\theta$, $\theta \in [0, \pi/2)$:
\[
[0,\tfrac\pi2) \ni \theta \longmapsto \zeta(x + i\tan\theta) \in \C \cup \{\infty\}.
\]
This maps each vertical half-line of the critical strip to a compact interval,
sending $i\infty$ to the point at infinity of the Riemann sphere $\hat\C$.
The image curve is a compactified trajectory of $\zeta$ that winds around $0$
once for each zero on $\Re(s) = x$.

The \emph{Volterra product integral} (product integral of the logarithmic derivative)
\[
F(x) := \Re\!\int_0^{\pi/2} \frac{\zeta'(x+i\tan\theta)}{\zeta(x+i\tan\theta)}\,
\frac{d\theta}{\cos^2\theta}
= \Re\!\int_x^{x+i\infty} \frac{\zeta'(s)}{\zeta(s)}\,ds
\]
counts zeros on the vertical line $\Re(s)= x$ via the argument principle.

\begin{conjecture}[Volterra detection of off-line zeros]\label{conj:volterra}
$F(x)$ is real-analytic on $(0,1)\setminus\{1/2\}$, with logarithmic
singularities precisely at values $x_0$ where $\zeta$ has a nontrivial zero on
$\Re(s) = x_0$. In particular, RH $\Leftrightarrow$ $F$ is smooth on
$(0,\tfrac12)\cup(\tfrac12,1)$.
\end{conjecture}

\begin{remark}[Why this stalls]
The integral $\int_x^{x+i\infty}(\zeta'/\zeta)\,dt$ is only conditionally
convergent. Proving $F(x_0) = 0$ for $x_0 \neq 1/2$ is equivalent to showing
that the oscillations of $\zeta'/\zeta$ along vertical lines cancel
perfectly --- which is essentially RH by another name. The geometric
compactification $t = \tan\theta$ clarifies the topology but does not remove
the analytic obstruction.
\end{remark}

\subsection{Infinite matrix traces and the von Neumann algebra}

View $\zeta(s) = \sum_{n=1}^\infty n^{-s}$ as the trace of the diagonal operator
$D^{(s)} = \mathrm{diag}(n^{-s})_{n\geq 1}$ on $\ell^2(\N)$:
\[
\zeta(s) = \mathrm{Tr}(D^{(s)}).
\]
The vertical product integral $F(x)$ is then the real part of the trace of
$\partial_t \log M^{(x,t)}$ integrated over $t \geq 0$:
\[
F(x) = -\Re\int_0^\infty \frac{d}{dt}\log\zeta(x+it)\,dt.
\]
A zero of $\zeta$ at $x+it_0$ produces a pole of $\partial_t\log M^{(x,t)}$ at
$t=t_0$, manifesting as a logarithmic divergence of the trace integral with
residue equal to the zero's multiplicity.

\subsection{Connection to Connes' trace formula}

Connes~\cite{Connes1999} constructs a spectral triple
$(\mathcal{A}, \mathcal{H}, \mathcal{D})$ with $\mathcal{H} = L^2(\A_\Q/\Q^*)$
and a self-adjoint Dirac-type operator $\mathcal{D}$. The spectral action is
\[
\mathrm{Tr}(f(\mathcal{D}/\Lambda))
\sim \Lambda\hat{f}(0)\zeta(1) - \sum_\rho \hat{f}(\rho) + O(1),
\]
where the sum runs over nontrivial zeros $\rho$.
The function $F(x)$ is a \emph{partial trace} of the resolvent of $\mathcal{D}$,
integrating out the imaginary direction.  RH is the statement that this partial
trace has no singular support off the single fiber $x = 1/2$.

The Dixmier trace $\mathrm{Tr}_\omega(T) = \lim_\omega \frac1{\log N}\sum_{n=1}^N \lambda_n(T)$,
applied to $\mathcal{D}^{-1}$, recovers $\zeta(1)$ and its distributional
prolongation. The vertical product integral at $x=1$ is this Dixmier trace,
regularized by $t = \tan\theta$.

\subsection{Heisenberg algebra over primes: integers as words, primes as letters}

Let $\mathcal{P} = \{2,3,5,7,\ldots\}$ be the primes (an alphabet).  Each
integer $n = \prod_p p^{v_p(n)}$ is a word in the free commutative monoid
$\mathrm{Mon}(\mathcal{P}) \cong (\Z_{>0}, \times)$.

Promote each prime $p$ to bosonic creation/annihilation operators
$a_p^\dagger, a_p$ on Fock space $\mathcal{F} = \bigoplus_{n\geq 0} \mathrm{Sym}^n(\ell^2(\mathcal{P}))$:
\[
[a_p, a_q^\dagger] = \delta_{pq}, \qquad
[a_p, a_q] = [a_p^\dagger, a_q^\dagger] = 0.
\]
Each integer $n$ becomes the Fock state
$|n\rangle = \prod_p (a_p^\dagger)^{v_p(n)}/\sqrt{v_p(n)!}\;|0\rangle$.
The unique factorization theorem is the Poincar\'e-Birkhoff-Witt theorem
for the commutative enveloping algebra: every basis state $|n\rangle$ is unique.

The \emph{logarithmic number operator}
\[
H = \sum_p \log(p)\, a_p^\dagger a_p
\]
satisfies $H|n\rangle = \log(n)|n\rangle$, and its partition function is exactly
$\zeta$:
\[
\mathrm{Tr}(e^{-\beta H}) = \sum_{n=1}^\infty n^{-\beta} = \zeta(\beta), \qquad \beta > 1.
\]

\begin{remark}[Primes as left-moving modes]
In a 2d CFT, the Virasoro character $\mathrm{Tr}_\mathcal{H}(q^{L_0})$ counts
operator states built from mode applications $L_{-n_1}\cdots L_{-n_k}|h\rangle$.
Here $\mathrm{Tr}(q^H) = \zeta(-\log q)$ counts integer states built from prime
modes $a_p^\dagger$. Primes are the simple roots of the commutative mode algebra;
unique factorization is PBW.
\end{remark}

\subsection{The center of the Heisenberg algebra and Hecke operators}

Operators commuting with all mode products $a_p^\dagger a_q$ are exactly the
\emph{multiplicative arithmetic functions} $f: \Z_{>0} \to \C$ (diagonal
operators $T = \sum_n f(n)|n\rangle\langle n|$ with $f(mn) = f(m)f(n)$ for
$\gcd(m,n)=1$). Their Dirichlet series $\sum f(n) n^{-s}$ are the Euler-product
$L$-functions.

The Hecke operators $\{T_p\}$ generate a commutative (Hecke) algebra whose center
consists of spherical functions --- automorphic $L$-functions for
$\mathrm{GL}(1)$ (Dirichlet characters) and $\mathrm{GL}(2)$ (classical modular forms).

\subsection{Connection to the Langlands program}

The four-floor algebraic tower is:
\begin{enumerate}
\item \textbf{Arithmetic} ($\Z$, $(\Z_{>0},\times)$): commutative, primes generate.
\item \textbf{Heisenberg} (CCR over $\mathcal{P}$): non-commutative in
  creation/annihilation; diagonal (abelian) in the number basis.
  Center = multiplicative arithmetic functions.
\item \textbf{Hecke algebra}: acts on Hilbert space of arithmetic functions;
  automorphic side of Langlands for $\mathrm{GL}(1)$.
\item \textbf{Virasoro / $W_\infty$}: modes of $T(z) = :(\partial\phi)^2:$
  built from the free boson on the charge lattice $(\log\mathcal{P})\subset\R$.
  Central charge $c = c_{\mathrm{eff}} = 2\alpha$ connects back to \S\ref{sec:virasoro}.
\end{enumerate}

Each floor is a non-abelian extension of the one below. RH on the ground floor
translates to: the spectrum of $L_0$ (Virasoro zero mode, fourth floor) has all
eigenvalues $h = h_{r,s}$ real and positive --- the Kac-table condition of
\S\ref{sec:virasoro}.

The Bost-Connes algebra~\cite{BostConnes1995}
$\mathcal{A}_{\mathrm{BC}}$ (with generators $\mu_n$, $e(r)$,
$r\in\Q/\Z$) provides the $\mathrm{C}^*$-algebra bridge between floors~2 and~3:
its KMS states at inverse temperature $\beta > 1$ give $\zeta(\beta)$ as
partition function, the same $\zeta$ appearing in the Fock-space trace at
level~2.

\begin{conjecture}[Center conjecture]\label{conj:center}
The center of the Virasoro algebra acting on $\mathcal{F}$ is generated by the
Casimir $C = L_0 - c/24$. Reality of the spectrum of $C$ is equivalent to
$c_{\mathrm{eff}} \in \R$, which is the real-grading condition of
\S\ref{sec:virasoro}, which implies RH.
\end{conjecture}

% ============================================================
\section{Structural Similarities: One Theorem in Four Languages}
\label{sec:similarities}
% ============================================================

All four approaches are the same meta-theorem in different languages.
We make the dictionary explicit.

\subsection{The common spectral operator and its involution}

Each approach defines an operator $\mathcal{O}$ whose spectrum (or whose kernel,
or whose resolvent singularities) encodes the zeros of $\zeta$:

\noindent
\textbf{Weyl derivative.} $\mathcal{O} = \partial^s$, a multiplicative operator
on Fourier modes. The critical line is the locus where the even/odd symmetry
$a_n = \pm a_{-n}$ holds, i.e., where the involution $f(\theta) \mapsto
f(-\theta)$ is compatible with $\partial^s f = 0$.

\noindent
\textbf{Schr\"odinger--Riemann.} $\mathcal{O} = -\tfrac12\partial_x^2$ acting on
$[0,\pi]$. The involution is $x \mapsto \pi - x$ (spatial reflection), which
enforces $A_n = A_{-n}$ and confines the ``forbidden normalization energies'' to
$\Re(s) = 1/2$.

\noindent
\textbf{Yang-Baxter / Virasoro.} $\mathcal{O} = L_0$ (Virasoro zero mode). The
involution is $s \mapsto 1-s$ (functional equation symmetry), encoded in the
Yang-Baxter system by the parabolic coordinate $w = s(s-1)$ which is
\emph{invariant} under $s \mapsto 1-s$. The real grading $\alpha \in \R$ is
the fixed-point condition under this involution.

\noindent
\textbf{Volterra / Langlands.} $\mathcal{O} = \mathcal{D}$ (Connes Dirac) or
$H = \sum_p \log p\,a_p^\dagger a_p$. The involution is adelic duality
$\A_\Q/\Q^* \to (\A_\Q/\Q^*)^\vee$, which corresponds to $n \mapsto 1/n$ on
the diagonal matrices, i.e., $s \mapsto 1-s$ on the Mellin transform.

\subsection{The symmetry condition in each language}

\begin{center}
\renewcommand{\arraystretch}{1.4}
\resizebox{\linewidth}{!}{%
\begin{tabular}{p{3.0cm} p{6.5cm}}
\textbf{Approach} & \textbf{Symmetry / reality condition} \\
\hline
Weyl derivative & $a_n = a_{-n}$: real (cosine) Fourier series \\
Schr\"odinger & $A_n = A_{-n}$: spatially symmetric wave function \\
Yang-Baxter & $\alpha \in \R$: real decay exponent / real grading of IBP system \\
Connes / Langlands & $F(x)$ real-analytic off $x=\tfrac12$; real spectrum of $L_0$ \\
\end{tabular}%
}\end{center}

In every case the condition is: some distinguished real-to-complex map
is real along the orbit of the natural involution $s \mapsto 1-s$.

\subsection{The Weyl derivative modes are Virasoro generators}

The Witt mode $d_n = -e^{in\theta}\partial_\theta$ is the $n$-th Fourier mode
of $\partial_\theta$, and $\partial^s = \sum_n (in)^s a_n e^{in\theta}$ is
the $s$-th complex power of $\partial_\theta$. On the Fock space $\mathcal{F}$
of \S\ref{sec:volterra}, $\partial_\theta$ corresponds to the generator $L_{-1}$
(translation), and $(in)^s$ corresponds to the spectral weight $n^s$ of the mode.
The Weyl decomposition \eqref{eq:weyl_decomp} is then:
\begin{align*}
\partial^s f &= \underbrace{i^s\,L(-s,\chi(n))}_{\text{left-moving / holomorphic}}
             + \underbrace{(-i)^s\, L(-s,\chi(-n))}_{\text{right-moving / anti-holomorphic}},
\end{align*}
the chiral decomposition of the fractional derivative into holomorphic
($a_n e^{in\theta}$) and anti-holomorphic ($a_{-n}e^{-in\theta}$) sectors.
Disjointness of $Z_+$ and $Z_-$ (Conjecture~\ref{conj:weyl}) is the condition
that the holomorphic and anti-holomorphic sectors have no common zero, which
is a version of the statement that the CFT has no operator at dimension $h = \bar{h}$
except on the diagonal --- reflecting the real structure of the operator algebra.

\subsection{Normalization and the partition function}

The Schr\"odinger normalization condition $\|\psi\|_{L^2} = 1$ and the
Yang-Baxter partition function $\mathrm{Tr}(e^{-\beta H}) = \zeta(\beta)$ are
the same object: both are the $L^2$ norm squared of a wave function
(or a density-matrix trace) in the respective Hilbert space.
The ``forbidden normalization energies'' of Conjecture~\ref{conj:qzeta}
are the zeros of $\zeta$; the ``spectral zeta function'' of the Schr\"odinger
Hamiltonian is $\zeta(s) = \sum_n E_n^{-s}$ by construction.

% ============================================================
\section{Generalizations of $\zeta$ via Seed and Quasi-Periodicity}
\label{sec:generalizations}
% ============================================================

The Yang-Baxter system has two free inputs: the seed integral and the
quasi-periodicity factor. Varying them generates a natural family generalizing $\zeta$.

\subsection{Varying the seed kernel}

Replace the Fermi-Dirac kernel $K(x) = e^x(e^x+1)^{-2}$ by any $K$ that is
smooth on $(0,\infty)$, rapidly decaying at $+\infty$, and controlled at $0$.
Representative examples:
\begin{itemize}
\item $K(x) = x^a e^{-x}$: seeds $\Gamma(s+a+1)$.
\item $K(x) = e^x(e^x-1)^{-2}$ (Bose kernel): same $q=1/4$ quasi-periodicity but
  different functional equation.
\item $K(x) = e^x(e^x+1)^{-n}$, $n \geq 2$: higher-order Fermi kernels,
  $q = 4^{-(n-1)}$.
\item $K(x) = J_\nu(x)e^{-x}$: Bessel kernels, Mellin transform related to
  $L$-functions of elliptic curves.
\end{itemize}

The seed determines the Euler product: $K(px)/K(x) = p^{-s}$ for all primes $p$
is the multiplicativity condition; its unique solution (up to normalization) is
the Fermi-Dirac kernel, explaining why $\zeta$ occupies a distinguished position.

In the Weyl framework, changing the seed corresponds to changing the Fourier
coefficients $a_n = \hat{K}(n)$: different kernels give different Dirichlet
series $L(-s, \chi(n))$, and the zero-set disjointness question becomes
theory-dependent.

\subsection{Varying the quasi-periodicity factor}

Replace $q = 1/4$ by $q \in (0,1)$. The quasi-periodicity becomes
$\sigma^{(q)}(s,n+2,m) = q\,\sigma^{(q)}(s,n,m)$. The Coulomb-gas identification
becomes $\alpha_+^{(q)}\alpha_-^{(q)} = q$, with the deformation parameter
$q_{\mathrm{YB}} = \sqrt{q}$ of the quantum group $U_q(\mathfrak{sl}(2))$.

\begin{conjecture}[Family of $L$-functions]\label{conj:qfamily}
For each $q \in (0,1)$ there exists an $L$-function $L^{(q)}(s)$ satisfying RH
if and only if the Yang-Baxter quadratic coefficients $a^{(q)}_{2j,2k}$ are real.
The family interpolates: $L^{(1/4)} = \zeta$, $L^{(q)} \to 1$ as $q \to 0$.
\end{conjecture}

% ============================================================
\section{Open Problems}
\label{sec:open}
% ============================================================

\begin{problem}[Algebraic determination of $c_{\rm eff}$]
Determine $c_{\mathrm{eff}} = 2\alpha \approx 2.220$ in closed form from the
Yang-Baxter recurrences.  See Problem~\ref{prob:ceff} for candidate approaches.
\end{problem}

\begin{problem}[Kac label map]
Find an explicit bijection $n \mapsto (r_n, s_n) \in \N^2$ such that
$t_n^2 = c_{\mathrm{eff}}|h_{r_n,s_n}| - 1/4$ holds for all nontrivial zeros.
\end{problem}

\begin{problem}[Proving disjointness of zero sets]
Prove (or disprove) Conjecture~\ref{conj:weyl}: for generic $\{a_n\}$,
$Z_+ \cap Z_- = \emptyset$. Even a proof in the case $a_n = n^{-\sigma}$
(standard Dirichlet series) would be interesting.
\end{problem}

\begin{problem}[Volterra conditional convergence]
Find a regularization of the vertical product integral $F(x)$ that
(a) is absolutely convergent and (b) still detects off-line zeros, without
presupposing RH.
\end{problem}

\begin{problem}[Langlands center condition]
Make Conjecture~\ref{conj:center} precise: identify the center of the Virasoro
algebra acting on $\mathcal{F}$ and show that real spectrum of $L_0$ implies
RH, without a circular argument.
\end{problem}

\begin{problem}[De~Sitter zeta analogue]
Compute $\sigma_{dS}(2,0,0)$ perturbatively in $T_{dS}/m$ and verify
Conjecture~\ref{conj:topology} quantitatively.
\end{problem}

% ============================================================
\section{Summary}
\label{sec:summary}
% ============================================================

\begin{center}
\renewcommand{\arraystretch}{1.4}
\resizebox{\linewidth}{!}{%
\begin{tabular}{p{2.4cm} p{3.6cm} p{3.4cm} p{2.4cm}}
\textbf{Approach} & \textbf{Main tool} & \textbf{Symmetry condition} & \textbf{Status} \\
\hline
Weyl derivative
  & $\partial^s = i^s L(-s,\chi(n)) + (-i)^s L(-s,\chi({-}n))$
  & $Z_+ \cap Z_- = \emptyset$ $\Rightarrow$ $a_n = a_{-n}$
  & Conjecture \\
Schr\"odinger
  & Forbidden normalization; $|D(s)|^2 = 0$
  & $L^2$ symmetry; $A_n = A_{-n}$
  & Heuristic \\
Yang-Baxter
  & Real grading; quadratic $as^2{-}as{+}1$
  & $\alpha \in \R$, $a < 4$
  & \textbf{Complete}~\cite{companion} \\
Volterra / Langlands
  & Arg.\ principle; Connes trace; Fock $\mathcal{F}$
  & Partial trace smooth off $x=\tfrac12$
  & Analytic gap \\
\end{tabular}%
}\end{center}

The Yang-Baxter quadratic-factor argument is the only approach that has been
closed into a rigorous proof.  The others remain as conjectures or heuristics,
each illuminating a different face of the same spectral reality condition that
forces zeros onto the critical line.

\begin{thebibliography}{9}

\bibitem{companion}
D.~Remmenga,
\textit{A Yang-Baxter Representation of the $\zeta$ Function},
preprint (2026).

\bibitem{paraboliccompanion}
D.~Remmenga,
\textit{Parabolic Resonance and a Formula for Riemann Zeros},
preprint (2026).

\bibitem{DotsenkoFateev1984}
V.~S.~Dotsenko and V.~A.~Fateev,
\textit{Conformal algebra and multipoint correlation functions in 2D statistical models},
Nucl.\ Phys.\ B \textbf{240}, 312 (1984).

\bibitem{Connes1999}
A.~Connes,
\textit{Trace formula in noncommutative geometry and the zeros of the Riemann zeta function},
Selecta Math.\ (N.S.) \textbf{5}, 29 (1999).

\bibitem{BostConnes1995}
J.-B.~Bost and A.~Connes,
\textit{Hecke algebras, type III factors and phase transitions with spontaneous symmetry breaking
in number theory},
Selecta Math.\ (N.S.) \textbf{1}, 411 (1995).

\bibitem{Hadamard1892}
J.~Hadamard,
\textit{Essai sur l'\'etude des fonctions donn\'ees par leur d\'eveloppement de Taylor},
J.\ Math.\ Pures Appl.\ \textbf{8}, 101 (1892).

\bibitem{Hermann2014}
R.~Hermann,
\textit{Fractional Calculus: An Introduction for Physicists},
World Scientific, Singapore, 2nd ed.\ (2014).

\end{thebibliography}

\end{document}
